\begin{table}[H]
\centering
\caption{Variable definitions and economic interpretations.}
\label{tab:variables}
\small
\begin{tabular}{@{}p{2.2cm}p{4cm}p{9cm}@{}}
\toprule
\textbf{Variable} & \textbf{Definition} & \textbf{Economic interpretation} \\
\midrule
GC contender & Equals 1 if the rider finished in the GC top 20 in any Grand Tour in the previous two years & Proxy for tournament incentives: GC contenders have higher expected returns from continued participation \\
Race: Tour de France & Dummy for the Tour de France (Vuelta as baseline) & Proxy for event value: the Tour has the highest media exposure, sponsor visibility, and UCI points \\
Race: Giro d'Italia & Dummy for the Giro d'Italia & Proxy for event value (intermediate prestige) \\
Previous DNF rate & Share of previous Grand Tour starts in which the rider abandoned & Proxy for rider-specific expected continuation cost or revealed risk type \\
Mountain stage & Equals 1 if the current stage is classified as a mountain stage (time-varying) & Proxy for marginal cost of effort: mountain stages impose the highest physical demands \\
Distance (z-score) & Standardised stage distance in km (time-varying) & Proxy for marginal cost of effort: longer stages are more physically taxing \\
Age & Rider's age at the start of the race & Human capital: captures physical maturity but also potential decline \\
Experience & Number of Grand Tours started prior to the current race & Human capital: captures accumulated knowledge and durability \\
Team size & Number of riders on the team starting the race & Team-resource control (ranges from 7 to 9) \\
Team average age & Mean age of the team's starting riders & Team-resource control \\
Year (centred) & Calendar year minus sample mean & Controls for secular trends in abandonment \\
\bottomrule
\end{tabular}
\end{table}